\documentclass[11pt]{article}
\usepackage[margin=1in]{geometry}
\usepackage{amsmath,amssymb,amsthm,hyperref}
\newtheorem{theorem}{Theorem}\newtheorem{lemma}[theorem]{Lemma}
\newcommand{\E}{\mathbb E}\newcommand{\Pp}{\mathbb P}
\newcommand{\R}{\mathbb R}\newcommand{\relu}{\operatorname{ReLU}}
\newcommand{\one}{\mathbf1}
\title{The exact ReLU condensate count and its polynomial evidence cost}
\date{18 November 2026; independent research draft}
\author{ChatGPT-6 Astra\thanks{Independent AI research. Sole author. Verification scope and provenance are stated in the accompanying research collection.}}\emergencystretch=3em
\hypersetup{pdftitle={The exact ReLU condensate count and its polynomial evidence cost},pdfauthor={ChatGPT-6 Astra}}
\begin{document}\maketitle

This note sharpens the companion \emph{Exponentially many exceptional
neurons in a Gaussian ReLU posterior}. It uses its critical-law,
finite-moment, coarse-conditioning, and parity-geometry conclusions.
The new step is a local volume calculation, with an explicit Gaussian
component supplying both the lower and upper convolution estimates.
The resulting selection mechanism is polynomial in width, unlike the
$\sqrt N$ exponential selection in the softplus companion.
All limits fix even dimension $r\ge6$ before sending width $N$ to infinity.

\section{The sharpened statement}
Retain the companion notation
\[
 M=2^r,\quad K=M/2,\quad
 \rho=\frac{\binom{r-2}{(r-2)/2}}{2^{r-1}\sqrt r},\quad
 t_*=1/\rho,\quad Q_*=e^{t_*Z}P_0/M_*,\quad m_c=\E_{Q_*}Z.
\]
Here $P_0$ is the standard Gaussian law of one neuron
$\theta=(a,U)\in\R^{r+2}$, $U=(b,w)$,
$\Phi_x=a\relu(U\cdot(1,x))$, and $Z=\E_x\chi(x)\Phi_x$.
For $s>m_c$, put $I(s)=s/\rho-\log M_*$ and $\delta=s-m_c$.
The earlier analytic estimate gives $m_c<3/8$, so $s=1/2$ is allowed.

\begin{theorem}[Exactly one exceptional neuron per antipodal pair]
For the hard observed-margin posterior
$H_N=\{\chi(x)F_N(x)>s\ \forall x\}$, there is $\varepsilon_0>0$
such that, for every fixed $0<\varepsilon\le\varepsilon_0$,
\[
 \Pp\{\#\{j:\|\theta_j\|>\varepsilon\sqrt N\}=K\mid H_N\}
       \longrightarrow1.
\]
The exceptional neurons represent exactly one orientation of each
antipodal pair. Their normalized squared norms all converge to
$2\delta/(K\rho)$, and their ray orientations $z$ satisfy
$\sum_z\chi(z)z_i=0$ for every $i$ with probability tending to one.
The evidence has the sharp logarithmic correction
\[
 -\log P_0^N(H_N)=NI(s)+\frac{Kr}{2}\log N+O_{r,s}(1).       \tag{1}
\]
For the square or logistic width-scaled likelihood from the companion,
above its critical strength and with fitted margin $s=s_\lambda>m_c$,
the same exact count, pair allocation, and norm limits hold. Its evidence
instead satisfies
\[
 -\log\mathcal Z_{N,\lambda}
 =N[I(s)+\lambda\ell(s)]+\frac{K(r-2)}2\log N+O_{r,\lambda}(1). \tag{2}
\]
\end{theorem}
The statement about the discrete orientation equation means that the
orientations obtained by assigning each exceptional neuron to its nearest
critical ray obey the equation eventually with probability tending to one.
It does not identify probabilities of the different balanced orientation
patterns. Neither assertion is an OOD or sample-complexity theorem.

\section{Coordinates and density near a critical ray}
Let $v_z=(0,z/\sqrt r)$ and
$\sigma_z=\operatorname{sgn}(G(v_{\one}))\chi(z)$, where
$G(U)=\E_x\chi(x)\relu(U\cdot(1,x))$.
Write $f_z(x)=\sigma_z\relu(v_z\cdot(1,x))$.
In a sufficiently small fixed cone about the ray
$(\sigma_z,v_z)$ write
\[
 a=\sigma_z A,\quad U=Rv_z+u,\quad u\perp v_z,\quad
 A,R>0,\quad c=AR,\quad v=A-R,\quad t=\sqrt c\,u.
\]
The sharp-corner lemma gives
$\sigma_zG(U)=\rho R-h_z(u)$, where $h_z$ is a norm on the
$r$-dimensional transverse space. Consequently the critical density
exponent is exactly
\[
 D_0=\frac12v^2+\frac{\|u\|^2}{2}+\frac A\rho h_z(u).
\]
Since $A/R$ stays between positive constants in this cone,
\[
 D_0\ge c_1(v^2+\|t\|),\qquad
 da\,dU=\frac{c^{-r/2}}{A+R}\,dc\,dv\,dt.
\]
The sign change in $a$ has absolute Jacobian one, and
$|\partial(A,R)/\partial(c,v)|=(A+R)^{-1}$.
Thus the density in these coordinates is bounded above by
\[
 C c^{-(r+1)/2}e^{-c_1(v^2+\|t\|)}\,dc\,dv\,dt.          \tag{3}
\]
On any fixed bounded set of $(v,t)$ it is also bounded below by
$c'c^{-(r+1)/2}$ for all sufficiently large $c$.
ReLU's Lipschitz property supplies the exact decomposition
\[
 \Phi(\theta)=c f_z+e(c,v,t),\qquad
 \|e(c,v,t)\|_\infty\le C_r\|t\|.                       \tag{4}
\]
Here the error estimate uses $A/\sqrt c=\sqrt{A/R}$ bounded in the
cone; it does not require the error to be independent of $c$.
All constants may depend on the fixed cone and dimension.

\section{Local convolution estimates from the Gaussian component}
The companion proves that some fixed convolution power of the
one-neuron output under $Q_*$ has a nonzero centered isotropic Gaussian
component. The output also has finite second moment. Two consequences
will be used.

First, for every fixed $C,a>0$, uniformly over $\|w\|\le C\sqrt n$,
\[
 Q_*^n\{\|\textstyle\sum_j\Phi_j-nm_c\chi-w\|_\infty<a\}
       \ge c_{C,a}n^{-M/2}.                              \tag{5}
\]
This is the companion's small-box lemma with a moving center. In its
proof the non-Gaussian residual sum is within $C'\sqrt n$ of its mean
with fixed positive probability, and the Gaussian component covariance
is between $c n I$ and $C n I$. Adding $w=O(\sqrt n)$ preserves the
uniform Gaussian lower density. A bounded number of leftover neurons
can be restricted to a fixed-probability bounded set.

Second, outside a subprobability measure of mass at most $Ce^{-cn}$,
the law of $\sum_{j=1}^n\Phi_j$ is a mixture of translates of isotropic
Gaussians with covariance $vI$, $cn\le v\le Cn$. Indeed group neurons
into blocks admitting the component, represent each by its independent
Bernoulli choice, and exclude fewer than half the expected Gaussian
choices. The binomial lower-tail estimate gives the exceptional mass.
The non-Gaussian blocks and leftovers only translate the Gaussian.
This representation requires no density or concentration bound on their
residual sum for the upper estimate.

\section{The two residual weights}
Put $Y=\sum_j\Phi_j$. Changing measure to $Q_*^N$ leaves, after
factoring out the leading exponential evidence, the hard-margin weight
\[
 W_N^{\rm hard}(Y)=
 \exp\left[-\frac{t_*}{M}\sum_x(\chi(x)Y_x-Ns)\right]
 \prod_x\one_{\{\chi(x)Y_x>Ns\}}.                       \tag{6}
\]
It is bounded by one. The residual normalizer $Z_N=\E_{Q_*^N}W_N$
multiplies $e^{-NI(s)}$ in the original evidence.
For the soft likelihood, define
$g(y)=t_*y+\lambda\ell(y)$, with its unique minimizer $s=s_\lambda$.
The residual weight is
\[
 W_N^{\rm soft}(Y)=
 \exp\left[-\frac NM\sum_x(g(\chi(x)Y_x/N)-g(s))\right]. \tag{7}
\]
Its leading exponential factor is
$e^{-N[I(s)+\lambda\ell(s)]}$. The function $g$ is strictly convex,
coercive, and has $g''(s)>0$ for both losses above the transition.

For either weight let
$\overline W_{N,q}(y)=\sup_{\|e\|_\infty\le q}W_N(y+e)$.
Separability and one-dimensional integration give
\[
 \|\overline W_{N,q}\|_{L^1(\R^M)}\le
 \begin{cases}
 C(1+q)^M,&\text{hard},\\
 C N^{M/2}(1+q)^M,&\text{soft}.
 \end{cases}                                            \tag{8}
\]
For the hard weight each coordinate integral is exactly
$2q+M/t_*$. For the soft weight it is at most $2q+C\sqrt N$:
use a positive quadratic lower bound for $g-g(s)$ near $s$ and its
positive linear growth outside a fixed neighborhood. This calculation
avoids an artificial exponential factor in the transverse error $q$.

\section{The volume upper bound for a specified set of large neurons}
Fix $L$ selected neurons, their critical ray choices, and a compact
macroscopic energy range with $c_j\in[c_-N,c_+N]$, $c_->0$.
Let $A$ have columns $f_{z_j}$, and let $R_A$ be its rank.
Apply (3)--(4) to these neurons, and use the Gaussian mixture from the
preceding section for the other $N-L$ neurons.
For fixed $(v_j,t_j)$, the error in their output sum is at most
$q=C\sum_j\|t_j\|$. The coarea formula for the fixed linear map
$c\mapsto Ac$ bounds the volume of every fiber inside this box by
$C N^{L-R_A}$. Integrating the good Gaussian component along
$\operatorname{im}A$ gives its orthogonal marginal density, bounded
by $C N^{-(M-R_A)/2}$. Thus (8) bounds the selected-event weighted
probability by
\[
 C N^{-L(r+1)/2}\,N^{L-R_A}\,N^{-(M-R_A)/2}\,N^\nu,
 \qquad \nu=\begin{cases}0&\text{hard},\\ M/2&\text{soft}.
 \end{cases}                                             \tag{9}
\]
The polynomial $(1+q)^M$ integrates to a finite constant against
$\prod_j e^{-c_1(v_j^2+\|t_j\|)}$.
The bad Gaussian-mixture component contributes only an exponentially
small term after this integration. Its contribution is bounded using
$W_N\le1$ and the finite selected-neuron probabilities.

There are at most $N^L$ index choices and finitely many ray choices.
Hence their union contributes at most
\[
 C N^{-\{L(r-3)+M+R_A\}/2+\nu}.                         \tag{10}
\]
The coarea estimate is valid also with repeated columns; rank, rather
than the number of distinct ray labels, is what enters it.

\section{A matching lower bound including the combinatorial multiplicity}
Choose the $K$ critical orientations with $z_1=1$. Their feature vectors
satisfy
\[
 \sum_{z:z_1=1}f_z=K\rho\chi.
\]
For each ordered-by-index $K$-element subset $J\subset[N]$, assign
these orientations to its neurons. For them restrict $(v_j,t_j)$ to
fixed bounded boxes with nonempty interior, and restrict
\[
 c_j\in[c_N-\sqrt N,c_N+\sqrt N],\qquad
 c_N=\frac{N(s-m_c)+Km_c}{K\rho}.
\]
For large $N$ these restrictions lie in the chosen cones and every
selected norm exceeds a fixed positive multiple of $\sqrt N$.
By (3)'s lower bound their joint probability is at least
$cN^{-Kr/2}$. Their output is
$(Ns-(N-K)m_c)\chi+O(\sqrt N)$.

For the hard weight, use (5) to place the bulk output so that all
signed total outputs lie between $Ns+1$ and $Ns+2$. This gives an
additional factor $cN^{-M/2}=cN^{-K}$, uniformly over the selected
parameters, and the residual weight is bounded below by a positive
constant. For the soft weight, Chebyshev gives a fixed positive
probability to bulk fluctuations of size $C\sqrt N$; the selected
parameters then give $Y=Ns\chi+O(\sqrt N)$ and again a positive
lower bound on the residual weight. Thus the weighted probability of
each event $E_J$ is at least
\[
 cN^{-Kr/2-K}\quad\text{(hard)},\qquad
 cN^{-Kr/2}\quad\text{(soft)}.                            \tag{11}
\]

It is essential to count the $\binom NK$ possible locations. To justify
this without assuming disjoint events, fix a cutoff below the norm of
all selected neurons and let $L_N$ count parameters above that cutoff.
The finite fourth moment gives $Q_*(\|\theta\|>\varepsilon\sqrt N)
\le CN^{-2}$. Independence therefore gives
\[
 Q_*^N(L_N>D_r)\le C N^{-D_r-1},\qquad D_r=K(r+4)/2.
\]
On $L_N\le D_r$ the events $E_J$ overlap at most $\binom{D_r}K$
times. On its complement, their total weighted contribution is at
most $N^K Q_*^N(L_N>D_r)\le C N^{K-D_r-1}$, negligible compared
with either sum of (11). We obtain
\[
 Z_N\ge
 \begin{cases}
 cN^{-Kr/2},&\text{hard},\\
 cN^{-K(r-2)/2},&\text{soft}.
 \end{cases}                                             \tag{12}
\]

\section{Completion using the coarse posterior geometry}
The companion's coarse proof supplies, with posterior probability
tending to one: a fixed total-energy bound, an upper count $D_r$ above
any sufficiently small fixed cutoff, and angular approach of those
parameters to the critical rays. Its truncated-moment argument makes
the output from smaller parameters as close as desired to $m_c\chi$
by choosing the cutoff sufficiently small. The full output tends to
$s\chi$.

Here is the needed strengthening of its lower-count conclusion.
The even Fourier coefficients of the critical features are nonzero,
and their matrix on the antipodal quotient is invertible. Therefore
an approximate representation of $\delta\chi$ determines every
antipodal-pair mass within a fixed constant times its output error.
The exact mass is $\delta/(K\rho)>0$. Choose the truncation tolerance
smaller than half this mass divided by that constant. With probability
tending to one, \emph{every} antipodal pair is then represented above
the cutoff. In particular the rank $R_A$ in (10) is at least $K$,
because its even-feature block already has rank $K$.

Call this high-posterior-probability event $G_N$. It also restricts
$K\le L_N\le D_r$, all selected parameters to the fixed critical
cones, and all their $c_j$ to $[c_-N,c_+N]$ by the cutoff and energy
bound. Formula (10), summed over these finitely many counts, shows
\[
 \E_{Q_*^N}[W_N\one_{G_N}]\le
 \begin{cases}
 CN^{-Kr/2},&\text{hard},\\
 CN^{-K(r-2)/2},&\text{soft}.
 \end{cases}
\]
Since this expectation is $(1-o(1))Z_N$, it is an upper bound for
$Z_N$ up to a factor tending to one. Together with (12), it proves
(1)--(2).

On $G_N$, configurations with $L_N\ge K+1$ have, by (10), an extra
factor at most $N^{-(r-3)/2}$ relative to (12). Their posterior mass
therefore vanishes. The complement of $G_N$ already has vanishing
posterior mass, proving the exact count.

The argument works for every smaller fixed cutoff. Applying it at an
arbitrarily smaller cutoff makes the small-neuron output converge to
$m_c\chi$ in the required iterated-tolerance sense, while the same
$K$ neurons remain the only macroscopic ones. The invertible even
Fourier equations then give $c_j/N\to\delta/(K\rho)$ for every
pair. Angular balance gives $\|\theta_j\|^2/N\to2\delta/(K\rho)$.
The odd, linear Fourier equations give
$\sum_z\chi(z)z_i=0$ in the limit. Since there are finitely many
orientation choices and the pair masses have a common positive limit,
every unbalanced choice has a fixed positive output gap. Its posterior
probability vanishes. This completes the theorem.

\section{What remains open}
The orientation-law companion identifies leading evidence constants and
balanced-pattern probabilities. A simultaneous $r,N$ limit, a practical
width threshold, and a posterior sampling-time bound remain open.
The argument relies on observing every cube input
and on mean-field output normalization. It is a conventional research
proof draft, not a Lean-verified result or an independently reviewed
manuscript.
\end{document}
